% ---------------------------------------------------------------------------
%  Embedded Spaces for Discovering New Classes
%  ICLR submission version (target 9-10pp main text).
%
%  STATUS: Sections 1-3 drafted 2026-08-26 (intro + algorithms).  Results
%  sections are DELIBERATELY not written -- exps 124/125/126/131 are pending on
%  the GPU machine and every number in them would be provisional.  Placeholders
%  below say exactly what each results section is waiting on.
%
%  The long-form technical report (discovery_metrics_iclr.tex, 41pp) stays as
%  the companion; this is the submission.  Preamble emulates the ICLR style so
%  it compiles with stock pdflatex; see that file for the same block.
% ---------------------------------------------------------------------------
\documentclass{article}

% ------------------------- ICLR-EMULATION (begin) -------------------------
\usepackage[letterpaper,textwidth=5.5in,textheight=9in,top=1in,left=1.5in]{geometry}
\usepackage{times}
\renewcommand{\ttdefault}{cmtt}
\renewcommand{\baselinestretch}{1.0}
\makeatletter
\renewcommand\section{\@startsection{section}{1}{\z@}%
  {-2.0ex \@plus -0.5ex \@minus -0.2ex}{1.2ex \@plus 0.2ex}%
  {\large\bf\raggedright}}
\renewcommand\subsection{\@startsection{subsection}{2}{\z@}%
  {-1.6ex \@plus -0.4ex \@minus -0.2ex}{0.9ex \@plus 0.2ex}%
  {\normalsize\bf\raggedright}}
\renewcommand\paragraph{\@startsection{paragraph}{4}{\z@}%
  {1.2ex \@plus 0.4ex \@minus 0.2ex}{-1em}{\normalsize\bf}}
\makeatother
% -------------------------- ICLR-EMULATION (end) --------------------------

\usepackage{amsmath,amssymb,amsfonts}
\usepackage{booktabs}
\usepackage{graphicx}
\usepackage{url}
\usepackage[colorlinks=true,linkcolor=black,citecolor=black,urlcolor=blue]{hyperref}
\usepackage{caption}
\captionsetup{font=small,labelfont=bf}

\newcommand{\E}{\mathbb{E}}
\newcommand{\R}{\mathcal{R}}
\newcommand{\D}{\mathcal{D}}
\newcommand{\pool}{\mathcal{P}}
\newcommand{\nmin}{n_{\min}}

\title{Embedded Spaces for Discovering New Classes}
\author{P.~C.~Harris\\ Massachusetts Institute of Technology \\ \texttt{pcharris@mit.edu}}
\date{}

\begin{document}
\maketitle

\begin{abstract}
% TODO: write last, once Sections 4-6 have numbers.
\end{abstract}

% ===========================================================================
\section{Introduction}
% ===========================================================================

A representation can encode a novel class perfectly and still be useless for
finding it. This is not a hypothetical failure mode; it is the ordinary
situation. Train a supervised contrastive encoder on a set of known classes,
hold one class out, and a linear probe trained to separate the held-out class
from the rest reaches an AUC of $0.94$. The same space, asked to flag novel
examples without being told what to look for, scores exactly zero at the
event level. The information is present and linearly accessible. It is simply
not accessible to anyone who does not already have the answer.

The gap between those two numbers is the subject of this paper. It is the gap
between a representation being \emph{decodable} and being
\emph{discoverable}, and it matters because scientific discovery lives
entirely on the second side of it. A search for a new particle, a new galaxy
morphology, or a new failure mode of a deployed model does not come with
labels for the thing being searched for. What it has is a large corpus, a
smaller reference sample of what is already understood, and a question: is
there something here that the reference does not explain?

Most of the representation-learning literature optimises for the first side.
Linear probe accuracy, $k$-NN accuracy, and transfer benchmarks all measure
how well a space supports a decision that someone else has already specified.
They are the right metrics for that purpose. They are close to uninformative
about whether a space will surface structure nobody has named yet, and --- as
the two numbers above indicate --- they can be actively misleading, because
the training objectives that maximise them are not the objectives that make
distances mean anything in absolute terms.

\paragraph{What we do.}
We treat discovery as the target and work backwards. The pipeline we study is
the standard open-world loop: score every point in a corpus against a
reference of labelled examples, take the most anomalous ones as a
\emph{pool}, cluster the pool to propose new class anchors, and adapt the
representation to the proposal. Every stage of that loop is downstream of one
quantity --- the fraction of the pool that is genuinely novel, which we call
\emph{pool purity}. Purity is the only forward-looking metric in the
pipeline. Every other measure describes a space you already have; purity
predicts whether the next step will improve it or corrupt it. Below a
threshold the loop does not merely fail to help: it clusters background,
plants anchors on it, and fine-tunes the representation to make that
structure sharper.

\paragraph{Contributions.}
\begin{enumerate}
\item \textbf{An algebra for the pool, and a ceiling.} Writing $b$ for the
  novel base rate, $q$ for the pool fraction, and $\varepsilon_s$ for the
  fraction of novel points retained, pool purity is
  $\Pi = \varepsilon_s b / q$, so the enrichment $\Pi/b$ is exactly
  $\varepsilon_s/q$ and
  $\Pi \le \min(1, b/q)$. The ceiling is a property of the cut alone. It says
  that a pool threshold inherited without derivation --- as ours was --- can
  cap achievable purity far below the gate no matter how good the
  representation is, and it separates ``the space is bad'' from ``the cut is
  bad'' for the first time.
\item \textbf{A label-free rule for the cut.} Because purity requires the
  novel labels, tuning the cut on purity is oracle tuning. We derive the
  novel-count estimator that the Neyman--Pearson critic already provides, and
  select the tightest cut that retains enough \emph{estimated} novel points to
  remain clusterable. The rule uses only the corpus and the labelled
  reference. On synthetic data with an exact critic it recovers the oracle cut
  exactly; on real embeddings it raises pool purity from $0.232$ to $0.778$
  where the scorer is strong, and refuses --- reporting that it has refused
  --- where it is not.
\item \textbf{A regime split, stated as a condition rather than a caveat.}
  Single-class novelty at a rate of $\sim$1\% and multi-class novelty at
  $\sim$10\% are different problems with different answers: which pooling
  statistic wins, whether stricter thresholds help or hurt, and whether the
  gate is reachable at all invert between them. We report them separately
  throughout.
\item \textbf{A reduced set of objectives.} From a design cube of ten
  contrastive objectives we identify the few that earn a place for discovery
  specifically, and show that the ranking by discovery is not the ranking by
  probe.
\end{enumerate}

\paragraph{What this paper is not.}
We do not claim that a regulariser which makes embedded distance behave like a
log-likelihood thereby improves discovery; we tested that directly and it does
not hold (Section~\ref{sec:limits}). We do not claim the pipeline works
everywhere. A substantial part of what follows is a characterisation of where
it does not, and why --- in several cases the obstruction is the novel-class
rate rather than anything about the representation, which is a limit no
encoder can move.

% ===========================================================================
\section{Setup: discovery as a two-sample problem}
% ===========================================================================

\paragraph{The open-world protocol.}
We are given a corpus $\D$ of $N$ embedded points, of which an unknown
fraction $b$ belong to classes never seen in training, and a reference
$\R \subset \D$ of points carrying known class labels. No label for any novel
class is available at any point --- not for choosing hyperparameters, not for
setting thresholds. We state this explicitly because the constraint is easy to
violate accidentally, and because a method that quietly violates it reports
numbers that cannot be reproduced in the situation it claims to address.

\paragraph{The loop.}
One round proceeds as follows. Score every point by an anomaly score $s(x)$;
threshold to obtain a pool $\pool$; cluster $\pool$ and append the resulting
centroids to the anchor set; pseudo-label every pooled point by its nearest
discovered anchor; and adapt the representation (or, in the frozen variant of
Section~\ref{sec:constructions}, only the anchors) to the enlarged anchor set.
The round repeats with the updated anchors.

\paragraph{Pool purity, and why it gates everything.}
Let $\Pi$ denote the fraction of $\pool$ that is truly novel. The pool feeds
directly into a clustering step with no supervision downstream of it, so
$\pool$ is the last point at which a mistake can still be corrected and
nothing after it checks the work. At $\Pi \approx 0.01$ the clustering is not
finding the novel class; it is partitioning background, and the subsequent
fine-tune sharpens that partition. Empirically the probe drops by $0.031$ to
$0.038$ wherever such pools form. Purity therefore behaves as a gate rather
than a gradient: above roughly $0.15$ the loop compounds across rounds, and
below it the loop is harmful. We report $\Pi$ as the primary metric and treat
the gate as a precondition for running the pipeline at all.

% ===========================================================================
\section{Algorithms}
\label{sec:algorithms}
% ===========================================================================

This section describes the objectives that produce the embedding, the
statistic that scores it, and the rule that turns that statistic into a pool.
The three are usually treated as separate literatures; a contribution of this
work is that the middle one --- the Neyman--Pearson density ratio --- is the
same object in all three roles, which is what makes the last one derivable
rather than tuned.

\subsection{Objectives}

We work within a design cube spanning
$\{$instance, supervised$\}$ positives $\times$
$\{$cosine, bilinear, distance$\}$ critics $\times$
$\{$softmax, NPLM$\}$ estimators $\times$
$\{$none, SIGReg, class-wise SIGReg$\}$ marginals, which contains SimCLR,
SupCon, LeJEPA/VISReg and the density-ratio objectives as corners. Of the ten
arms we ran, three earn a place for discovery.

\paragraph{SupCon\,+\,SIGReg (strong).}
Supervised positives, cosine critic, softmax estimator, and a sliced
isotropic-Gaussian regulariser at weight $\lambda = 5$. This is the discovery
workhorse: it produces the highest pool purities we measure. Two points of
care. First, $\lambda$ matters and is not a formality --- the same objective
at the library default $\lambda = 1$ is a different algorithm and its numbers
must not be pooled with these (Appendix~\ref{sec:naming}). Second, we make no
claim that SIGReg helps \emph{because} it renders distance a log-likelihood;
that mechanism is tested and rejected in Section~\ref{sec:limits}. SIGReg
earns its place empirically, and structurally: under an isotropic marginal
Euclidean distance is Mahalanobis distance, which is what licenses the
unit-variance clustering criterion used downstream.

\paragraph{Supervised distance-NPLM.}
Supervised positives, distance critic, and the NPLM estimator
\begin{equation}
\label{eq:nplm}
\mathcal{L}[g] \;=\; \E_{\R}\!\left[e^{g} - 1\right] \;-\; \E_{+}\!\left[g\right],
\end{equation}
whose minimiser is the \emph{absolute} log density ratio, fixed by the
calibration $\E_{\R}[e^{g^\star}] = 1$ with no free additive constant. Unlike
a softmax-normalised critic, this makes a single embedded distance a statement
about relative likelihood rather than a statement only about rank. It is the
strongest arm in the low-class-count regime and the most robust across
pretraining backbones.

\paragraph{SupCon (plain).}
Supervised positives, cosine critic, softmax estimator, no marginal. Retained
as the baseline and as the parent for the residual construction. It is the
best space in the paper by linear probe and by semantic agreement, and one of
the weakest by every calibration measure --- which is the dissociation this
paper is organised around, in a single arm.

We drop plain SimCLR, LeJEPA/SIGReg-SSL, instance NPLM-bilinear and instance
NPLM-distance to a single ablation row: none reaches the Pareto front for
discovery, and the bilinear critic in particular carries roughly three times
the seed variance of the distance critic.

\subsection{The Neyman--Pearson statistic, in three roles}
\label{sec:np}

Let $\R$ be the reference and $\D$ the corpus, with $w = N_\D / N_\R$. The
empirical NP objective
\begin{equation}
\label{eq:np}
\mathcal{L}[f] \;=\; w \sum_{x \in \R}\!\left(e^{f(x)} - 1\right) \;-\; \sum_{x \in \D} f(x)
\end{equation}
has minimiser $f^\star(x) = \log\!\big(p_\D(x)/p_\R(x)\big)$, obtained by
setting the functional derivative $p_\R e^{f} - p_\D$ to zero, and therefore
satisfies
\begin{equation}
\label{eq:calib}
\E_{\R}\!\left[e^{f^\star}\right] = \int p_\R \cdot \frac{p_\D}{p_\R} = 1 .
\end{equation}
Equation~\ref{eq:np} is Equation~\ref{eq:nplm} with pairs replaced by events.
The same functional appears as a contrastive \emph{loss} over pairs, as a
two-sample \emph{test statistic} $t_{NP} = -2\mathcal{L}$ distributed as a
Wilks log-likelihood ratio, and --- below --- as a \emph{pool scorer}. We
parameterise $f$ as a small kernel ensemble with learnable centres and
amplitudes and anneal the bandwidth from the median pairwise distance
downwards; a fixed bandwidth is calibrated only for one embedding scale and
goes blind when the geometry is rescaled.

Equation~\ref{eq:calib} is also a free diagnostic, and we report it. A fitted
critic need not satisfy it, and a value far from $1$ means the numbers built
from that fit are unreliable. Measured on trained spaces it is $0.997$ to
$1.019$ in-sample; out-of-sample it ranges to $2.7$, reflecting mild
overfitting of the reference subsample.

\paragraph{Why the ratio and not the distance.}
Distance-based pooling asks which points are far from everything. The density
ratio asks which points are over-represented relative to what the reference
explains. These come apart, and the difference is decisive: a novel class is
not a scatter of far-flung outliers but a \emph{dense clump} sitting where the
reference has no mass. On a $100$-class benchmark the extreme distance tail is
owned by background outliers, and tightening the distance threshold makes
purity \emph{worse} ($0.036 \to 0.002$), whereas tightening the ratio
threshold makes it better ($0.232 \to 0.358$). Same space, same rate, opposite
response.

\subsection{Choosing the pool: an algebra and a rule}
\label{sec:cut}

\paragraph{The algebra.}
Let $q = |\pool|/N$ and let $\varepsilon_s, \varepsilon_b$ be the fractions of
novel and background points retained. Then
$q = \varepsilon_s b + \varepsilon_b (1-b)$ and
\begin{equation}
\label{eq:purity}
\Pi \;=\; \frac{\varepsilon_s\, b}{q},
\qquad
E \;\equiv\; \frac{\Pi}{b} \;=\; \frac{\varepsilon_s}{q},
\qquad
\Pi \;\le\; \min\!\Big(1,\ \frac{b}{q}\Big).
\end{equation}
Three consequences. The enrichment is exactly the lift at operating point $q$.
Purity increases monotonically as $q$ shrinks, so there is no interior optimum
and targeting a purity \emph{level} is the wrong objective. And the ceiling
$\min(1, b/q)$ is a property of the cut alone: at $q = 0.05$ and $b = 0.01$ no
scorer, however good, can exceed $\Pi = 0.20$. A threshold chosen without
reference to $b$ can therefore cap purity near the gate independently of the
representation --- and the widely used ``$95$th percentile of the reference''
is such a threshold.

\paragraph{The rule.}
Purity cannot be used to choose $q$, since computing it requires the novel
labels. The critic supplies the missing quantity. Under the contamination
model $p_\D = (1-b) p_\R + b\, p_S$, with $r = e^{f}$, the posterior that a
corpus point is novel is $1 - (1-b)/r$; dropping the unknown $b$ and taking
the positive part gives the label-free weight
\begin{equation}
\label{eq:weights}
w(x) \;=\; \left[\,1 - \frac{1}{r(x)}\,\right]_{+},
\qquad
\E_\D[w] \;=\; \mathrm{TV}(p_\D, p_\R) \;=\; b\,\mathrm{TV}(p_S, p_\R) \;\le\; b ,
\end{equation}
so $\sum_{x \in S} w(x)$ estimates how many novel points a subset $S$
contains, exactly when the novel class is disjoint from the reference and
biased \emph{downward} under overlap --- a safe direction, since it tightens
the cut. Since purity rises as $q$ falls, the binding constraint is not purity
but whether a cluster survives, and that is what we target: ordering the
corpus by $f$ and writing $\widehat{N}(k)$ for the null-corrected sum of
$w$ over the top $k$ points,
\begin{equation}
\label{eq:rule}
k^\star \;=\; \min\big\{\,k : \widehat{N}(k) \ge \nmin \,\big\},
\qquad
q^\star = k^\star / N ,
\qquad
k_{\max} = \Big\lfloor \widehat{N}(N)\,/\,\nmin \Big\rfloor .
\end{equation}
The same weights bound the clustering search range, replacing a bound that in
the original pipeline was set from the number of held-out classes --- a
quantity an open-world method does not have.

\paragraph{The null correction, which is not optional.}
Equation~\ref{eq:weights} is a positive part, so \emph{any} spread in a fitted
$f$ manufactures apparent novelty. On pure-noise scores with no novel class at
all, the raw sum returns $665$ ``novel'' points out of $8000$ --- enough for
the rule to fire on nothing and hand the loop a pool of pure background.
Subtracting the mean weight of the reference does not repair this, because
top-ranked null points carry far above-average weight and their tail survives
subtraction. We instead subtract a \emph{rank-matched} null: under the null
hypothesis the top $q$ of the corpus is distributed as the top $q$ of the
reference, so
$\widehat{N}(k) = \sum_{i \le k} w(x_{(i)}) - k \cdot \overline{w}_{\R}^{(q)}$,
with $\overline{w}_{\R}^{(q)}$ the mean weight of the top-$q$ fraction of the
reference. This cancels to zero under the null while preserving a real excess.
When $\widehat{N}(N) < \nmin$ the rule declines to pool and reports that it has
declined; that refusal is an output of the method, and we report it as one.

\paragraph{Constants.}
$\nmin$ and the clipping range on $q$ are set once, globally, from a sweep on
held-out embeddings rather than inherited: purity falls monotonically as
$\nmin$ grows, so it is set as small as the clustering will bear.

% ===========================================================================
\section{Constructions}
\label{sec:constructions}
% ===========================================================================
% TODO. Two subsections, both already established and not waiting on new runs,
% but written after the results tables so the numbers can be quoted inline:
%   (a) residual decomposition + concatenation;
%   (b) frozen-space discovery -- freeze the encoder, iterate only the anchors
%       against a density-ratio pool.  The gain lives in the anchors, not the
%       representation update; unfreezing costs probe and, with distance
%       pooling, collapses round 2.
%   NOTE: re-verify the frozen numbers first -- the BatchNorm fix means the
%   archived frozen CIFAR results no longer reproduce (exp 130/131).

% ===========================================================================
\section{Experiments}
% ===========================================================================
% WAITING ON: exps 124 (leakage-free scratch shortlist), 125 (single-holdout
% battery with draws), 126 (restricted multi-holdout), 131 (cifar10/cifar100/
% galaxy10 re-run under the new pool cut).  Do not draft prose around
% provisional numbers.
%
% Table plan, split by regime and never pooled:
%   T1 single-holdout, all datasets, mean +- sd ACROSS DRAWS (not seeds --
%      draw spread exceeds seed spread, so seed-only intervals understate).
%   T2 multi-holdout, label-rich datasets only.
%   T3 pool-cut ablation: inherited q=0.05 vs the rule of Sec 3.3, with the
%      ceiling min(1, b/q) shown alongside so the reader can see headroom.
%
% Dataset ordering is by distance from pretraining, not alphabetical:
% galaxy10 (never seen by the backbone) -> DTD (task-novel) -> cars/flowers/
% aircraft (fine-grained in-domain) -> CIFAR (from-scratch lineage; the
% hub-pretrained cells saw the held-out class and are an ablation only).

% ===========================================================================
\section{When it does not work}
\label{sec:limits}
% ===========================================================================
% TODO, but the content is already settled and should not be softened:
%   - The rate floor. At b ~ 1% the gate is out of reach for the pipeline as
%     built. Note honestly that the CEILING is 0.20 there, not 0.03, so this
%     is a signal-efficiency limit and a clustering-detectability limit, NOT
%     an information-theoretic one.
%   - SIGReg does not help by making distance a log-likelihood. Forcing unit
%     width is decorative; seen-class fidelity anti-correlates with novelty
%     calibration. State the defensible claim instead.
%   - The temperature basin is unreachable by any label-free or transductive
%     selection: a landscape finding, not a recipe.
%   - Repairing an identity is not the same as improving a space (five
%     independent confirmations).
%   - Optimising the representation THROUGH a fitted NP kernel destroys the
%     separation it measures: the NP loss transfers as a form, not as an
%     optimisation target.

\appendix
\section{Objective naming}
\label{sec:naming}
% TODO: reproduce the canonical mapping from docs/LOSSES.md. Three objectives
% collide under similar names at two regularizer strengths; their numbers must
% not be pooled.

\end{document}
